\section{Practical Use, Parameters and Comparison}\label{sec:apps}

This section places Kronecker-FRI in its intended setting, extends it to the opening of several polynomials at one point, recommends parameter sets, and compares it with existing multilinear polynomial commitment schemes.

\subsection{Where the scheme fits}\label{sec:apps:fit}

\paragraph{Sumcheck-based SNARKs.}
Multilinear polynomial commitments are the cryptographic component of SNARKs built from the sumcheck protocol~\cite{LFKN92}, such as Spartan~\cite{Set20} and HyperPlonk~\cite{CBBZ23}.
In these systems the prover commits to multilinear polynomials, typically given by their tables on $B_n$, and the sumcheck reduces the statement to claims $f_i(z) = v_i$ at one random point $z \in \F^n$ chosen by the verifier; claims at several points are reduced to one point by a further sumcheck~\cite{CBBZ23}.
This is exactly the setting of $\Pi_{\mathrm{KF}}$ and of its batched version below: polynomials committed over the base field $\Fq$ (\cref{lem:enc}(3)), and one evaluation point in the extension $\F$.
A table is converted to the coefficient form of $f$ by the M\"obius transform before the NTT; this costs less than $1/(3 \cdot 2^R)$ of the commitment's arithmetic (\cref{sec:scheme:enc}), and $1$--$3\%$ of the commitment time in the measurements of \cref{sec:exp:scaling}.

\paragraph{Reuse of univariate FRI.}
Kronecker-FRI is a univariate low-degree test in disguise: the prover commits to univariate polynomials $U_f$ and $A$ of degree less than $N$, the verifier forms the batched word $w + \beta w_A + \beta^2 h$ pointwise, and the rest is FRI with the kernel fold (or, equivalently, the classical fold, \cref{rem:cfold}).
The virtual word $h$ is a pointwise quotient of the same kind as those used in FRI-based STARKs~\cite{BCRSVW19}.
A code base that already has an FRI prover and verifier over a smooth multiplicative domain therefore obtains a multilinear commitment by adding the product $U_f K_z$ on the prover side and the $O(n)$ computation of $h(\pm\xi)$ per query on the verifier side, with no sumcheck and no new proximity test.

\paragraph{Fields.}
The scheme needs a smooth domain of order $M = 2^{n+R}$ in $\Fq^\times$, that is, $2^{n+R} \mid q - 1$ (\cref{lem:power}).
Besides Goldilocks ($2^{32} \mid q-1$), this holds for the $31$-bit primes BabyBear, $q = 15 \cdot 2^{27} + 1$, and KoalaBear, $q = 127 \cdot 2^{24} + 1$, for $n + R \le 27$ and $n + R \le 24$ respectively; the extension degree $e$ must then make the field term $3M/q^e$ negligible, which requires $e \ge 5$ for $100$ bits at $M = 2^{24}$.
It does not hold for the Mersenne prime $2^{31}-1$, for which circle-domain analogues of FRI exist~\cite{HLP24}, nor in characteristic $2$ (\cref{sec:apps:compare}).
Only Goldilocks is implemented.

\subsection{Opening several polynomials at one point}\label{sec:apps:batch}

A SNARK prover typically opens several committed polynomials at the same point.
Opening them one by one costs one run of $\Pi_{\mathrm{KF}}$ each; the standard random linear combination reduces this to one run.
We state and prove it for the interactive protocol.

\begin{definition}[Batched evaluation protocol $\Pi^{m}_{\mathrm{KF}}$]\label{def:batch}
Let $m \ge 2$. The instance is
\[
  \bigl(w^{(1)}, \dots, w^{(m)};\ z,\ v_1, \dots, v_m\bigr),
\]
with oracles $w^{(i)} \in \F^L$, a point $z \in \F^n$ and values $v_i \in \F$; honestly $w^{(i)} = \Enc(f_i)$ and $v_i = f_i(z)$.
The protocol has $\ell + 3$ rounds.
In round $1$ the prover sends nothing and the verifier sends $\gamma \in \F$; put
\[
  w_\gamma \;=\; \sum_{i=1}^{m} \gamma^{i-1} w^{(i)}, \qquad v_\gamma \;=\; \sum_{i=1}^{m} \gamma^{i-1} v_i .
\]
Rounds $2$ to $\ell+3$ are rounds $1$ to $\ell+2$ of $\Pi_{\mathrm{KF}}$ on the instance $(w_\gamma; z, v_\gamma)$, where the verifier computes each value $w_\gamma(\pm\xi)$ from the $2m$ queries $w^{(i)}(\pm\xi)$.
The honest prover runs the honest prover of $\Pi_{\mathrm{KF}}$ for $f_\gamma = \sum_i \gamma^{i-1} f_i$.
\end{definition}

Completeness follows from \cref{thm:complete}, since $w_\gamma = \Enc(f_\gamma)$ and $v_\gamma = f_\gamma(z)$ by linearity of $\Enc$ and of evaluation.
The relation asks that the commitments be \emph{jointly} close to codewords, on one common set of fibres:
\begin{multline*}
  \mathcal{R}^{\delta,m}_{\mathrm{KF}} \;=\; \Bigl\{\, \bigl(((w^{(i)})_i; z, (v_i)_i), (f_i)_i\bigr) \;:\; \exists\, S \subseteq L \text{ fibre-closed with } |S| \ge (1-\delta)M \\
  \text{such that } w^{(i)} = \Enc(f_i) \text{ on } S \text{ and } f_i(z) = v_i \text{ for all } i \,\Bigr\}.
\end{multline*}
A witness satisfies $\Delta^{\mathrm{fib}}_0(w^{(i)}, \Enc(f_i)) \le \delta$ for every $i$, so it is unique when it exists (\cref{lem:enc}(2)); for $m = 1$ the relation is $\mathcal{R}^\delta_{\mathrm{KF}}$.

\begin{theorem}[Soundness of batched openings]\label{thm:batch}
Let $0 < \delta \le \delta^*$. If $\bigl((w^{(i)})_i; z, (v_i)_i\bigr)$ has no witness in $\mathcal{R}^{\delta,m}_{\mathrm{KF}}$, then every prover makes the verifier of $\Pi^m_{\mathrm{KF}}$ accept with probability at most
\[
  \frac{(m-1)M}{|\F|} + \varepsilon_{\mathrm{KF}}.
\]
\end{theorem}

\begin{proof}
By \cref{lem:randomised} it suffices to treat a deterministic prover. Its messages from round $2$ on are functions of $\gamma$ and of the later challenges, so for each fixed $\gamma$ rounds $2$ to $\ell+3$ are an execution of $\Pi_{\mathrm{KF}}$ on $(w_\gamma; z, v_\gamma)$ by a deterministic prover.
Let $\Gamma$ be the set of $\gamma \in \F$ such that $(w_\gamma; z, v_\gamma)$ has a witness in $\mathcal{R}^\delta_{\mathrm{KF}}$ (\cref{def:relation}).
For $\gamma \notin \Gamma$, \cref{thm:sound} bounds the acceptance probability by $\varepsilon_{\mathrm{KF}}$; it remains to show $|\Gamma| \le (m-1)M$.

Suppose $|\Gamma| > (m-1)M$. For $\gamma \in \Gamma$, let $g_\gamma$ be the witness; then $\Delta(w_\gamma, \Enc(g_\gamma)) \le \Delta^{\mathrm{fib}}_0(w_\gamma, \Enc(g_\gamma)) \le \delta$ (\cref{lem:fibre-dist}(1)), so $\Delta(w_\gamma, \Cc_0) \le \delta$.

\emph{Correlated agreement.} The words $w_\gamma$ form the curve generated by $(w^{(1)}, \dots, w^{(m)})$, of degree $t = m-1 \ge 1$, and more than $tM$ of its members are within distance $\delta$ of $\Cc_0$. By \cref{thm:curves}, there are codewords $\hat c_1, \dots, \hat c_m \in \Cc_0$ and a set of at least $(1-\delta)M$ points on which $w^{(i)} = \hat c_i$ for every $i$. Let $\mathrm{Agr}$ be the set of all points where these $m$ equalities hold, and $\hat f_i \in \Ll_n(\F)$ with $\Enc(\hat f_i) = \hat c_i$ (\cref{lem:enc}(1)).

\emph{The combination is the decoded codeword.} Let $\gamma \in \Gamma$. The codeword $\sum_i \gamma^{i-1}\hat c_i = \Enc(\sum_i \gamma^{i-1}\hat f_i)$ agrees with $w_\gamma$ on $\mathrm{Agr}$, so it is within distance $\delta$ of $w_\gamma$, and so is $\Enc(g_\gamma)$. By \cref{lem:rs-params}(3) they are equal, and by \cref{lem:enc}(1), $g_\gamma = \sum_i \gamma^{i-1} \hat f_i$.

\emph{The values.} Hence $\sum_{i=1}^{m} \gamma^{i-1}(\hat f_i(z) - v_i) = g_\gamma(z) - v_\gamma = 0$ for every $\gamma \in \Gamma$.
The polynomial $\sum_{i=1}^{m} Z^{i-1}(\hat f_i(z) - v_i) \in \F[Z]$ has degree at most $m-1$ and more than $(m-1)M \ge m-1$ roots, so it is zero (\cref{lem:roots}), and $\hat f_i(z) = v_i$ for all $i$.

\emph{A common set of fibres.} For $\xi \in L \setminus \mathrm{Agr}$, the polynomial $d_\xi(Z) = \sum_{i=1}^{m} Z^{i-1}(w^{(i)} - \hat c_i)(\xi) \in \F[Z]$ is nonzero, of degree at most $m-1$; let $\mathrm{Bad}$ be the set of all its roots, over all such $\xi$. Then $|\mathrm{Bad}| \le (m-1)\,|L \setminus \mathrm{Agr}| \le (m-1)\delta M < |\Gamma|$, so there is $\gamma \in \Gamma \setminus \mathrm{Bad}$.
For this $\gamma$, $w_\gamma$ and $\Enc(g_\gamma) = \sum_i \gamma^{i-1}\hat c_i$ agree exactly on $\mathrm{Agr}$: at $\xi \notin \mathrm{Agr}$ they differ by $d_\xi(\gamma) \ne 0$.
Since $\Delta^{\mathrm{fib}}_0(w_\gamma, \Enc(g_\gamma)) \le \delta$, they agree on a fibre-closed set $S$ with $|S| \ge (1-\delta)M$ (\cref{lem:fibre-dist}(2)), and $S \subseteq \mathrm{Agr}$. So every $w^{(i)}$ agrees with $\Enc(\hat f_i)$ on $S$, and $(\hat f_i)_i$ is a witness, a contradiction.
\end{proof}

\begin{corollary}[Knowledge and binding for batched openings]\label{cor:batch}
Let $\mathsf{Ext}^m$ output, for each $i$, the $\hat f_i \in \Ll_n(\F)$ such that $\Enc(\hat f_i)$ is the codeword of $\Cc_0$ within fibre distance $\delta^*$ of $w^{(i)}$, or $0$ if there is none; it runs in time polynomial in $mM$ (\cref{rem:decode}).
If a prover makes the verifier of $\Pi^m_{\mathrm{KF}}$ accept with probability greater than $(m-1)M/|\F| + \varepsilon_{\mathrm{KF}}$, then $(\hat f_i)_i$ is a witness. In particular $\Pi^m_{\mathrm{KF}}$ is evaluation binding for each of the $m$ commitments, with the same error.
\end{corollary}

\begin{proof}
By \cref{thm:batch}, a witness $(f_i)_i$ exists. Each $\Enc(f_i)$ is within fibre distance $\delta \le \delta^*$ of $w^{(i)}$, so it is the codeword found by $\mathsf{Ext}^m$ (\cref{lem:unique}), and $\hat f_i = f_i$. Binding follows as in \cref{lem:rbr-binding}, since the witness is determined by the commitments.
\end{proof}

\begin{remark}[Costs and scope of batching]\label{rem:batch}
Compared with one run of $\Pi_{\mathrm{KF}}$, the batched protocol has one more round (the challenge $\gamma$); the prover computes $f_\gamma$ ($mN$ operations) and $w_\gamma$ ($mM$ operations), and the verifier makes $2m$ instead of $2$ queries to the commitments per query point, that is, $m$ authentication paths instead of one in the compiled scheme; all other costs are those of a single opening.
We state the batched protocol for the interactive setting. Its relation is fibre-based, like that of \cref{sec:pq}, so the erasure analysis of \cref{sec:pq:kstate} is the natural route to a post-quantum version; we have not carried it out. The batched protocol is also not implemented in the current code.
\end{remark}

\subsection{Recommended parameters}\label{sec:apps:params}

\begin{table}[t]
\centering
\small
\setlength{\tabcolsep}{4pt}
\begin{tabular}{l c c r r rrrr}
\toprule
target & $\F$ & $\rho$ & $\kappa$ & field term & commit (s) & open (s) & verify (ms) & proof (KiB) \\
\midrule
$100$ bits & $\mathbb{F}_{q^2}$ & $1/4$ & 148 & $2^{-104.4}$ & 1.28 & 3.97 & 6.8 & 1080 \\
$100$ bits & $\mathbb{F}_{q^2}$ & $1/8$ & 121 & $2^{-103.4}$ & 2.63 & 7.07 & 5.5 & 933 \\
$128$ bits & $\mathbb{F}_{q^4}$ & $1/4$ & 189 & $2^{-232.4}$ & 1.26 & 6.74 & 8.5 & 1453 \\
post-quantum & $\mathbb{F}_{q^4}$ & $1/4$ & 248 & $2^{-232.4}$ & 1.27 & 6.73 & 11.4 & 1904 \\
\bottomrule
\end{tabular}
\caption{Recommended parameter sets, with measurements at $n = 20$, $\ell = n - 8$ (final polynomial of $256$ coefficients), salted trees, on the machine of \cref{sec:exp}. The field term is the bound $3M/|\F|$ of \cref{sec:sound:params}; the query term $(1-\delta^*)^\kappa$ is below $2^{-100}$ in the first two rows, $2^{-128}$ in the third and $2^{-168}$ in the fourth. The last row is a post-quantum level in the sense of \cref{cor:pq}, not a bit-security level comparable to the others.}
\label{tab:apps-params}
\end{table}

\Cref{tab:apps-params} gives four parameter sets, with $\delta = \delta^* = (1-\rho)/2$ and $\kappa$ from \cref{tab:kappa}.
\begin{itemize}
  \item \emph{$100$ bits, $\rho = 1/4$} is the default: both terms of $\varepsilon_{\mathrm{KF}}$ are below $2^{-100}$, so $\varepsilon_{\mathrm{KF}} < 2^{-99}$, for $M \le 2^{26}$, that is, $n \le 24$.
  \item \emph{$100$ bits, $\rho = 1/8$} trades a prover $1.8$--$2.1$ times slower for proofs $14\%$ smaller and a faster verifier; it is preferable when proofs are stored or sent many times.
  \item \emph{$128$ bits} needs a larger field, since the field term over $\mathbb{F}_{q^2}$ is only about $2^{-104}$ at this size; with $\mathbb{F}_{q^4}$ it is negligible, and $\kappa = 189$.
  \item \emph{Post-quantum} is the set of \cref{rem:pq-params}, for which \cref{cor:pq} gives $\varepsilon_{\mathrm{ext}}(2^{64}) < 2^{-31.8}$ against quantum adversaries making $2^{64}$ oracle queries, for one opening, for the compiler of~\cite{CDHZ25} (\cref{rem:scope}).
\end{itemize}
The first three rows are levels of the interactive protocol (\cref{thm:sound}). For the compiled scheme, the loss of the Fiat--Shamir transformation must be accounted for separately (\cref{rem:classical}).
We choose $\ell = n - 8$ because it is among the fastest settings for verification at $n = 20$ while reducing proof size by $19\%$ with respect to $\ell = n$ (\cref{tab:exp-stop}); proofs shrink further with smaller $\ell$, at the cost of $O(\kappa N_\ell)$ verifier operations.

\subsection{Comparison with existing schemes}\label{sec:apps:compare}

\begin{table}[t]
\centering
\footnotesize
\setlength{\tabcolsep}{2.5pt}
\begin{tabular}{l l l l l}
\toprule
scheme & setup, assumption & proof & opening prover & evaluation via \\
\midrule
Gemini~\cite{BCHO22} & SRS, pairings & $O(\log N)$ & $O(N)$ group ops & univariate KZG \\
Zeromorph~\cite{KT24} & SRS, pairings & $O(\log N)$ & $O(N)$ group ops & univariate KZG \\
Mercury~\cite{EG25} & SRS, pairings & $O(1)$ & $O(N)$ field, $O(N)$ scalar mults & univariate KZG \\
Samaritan~\cite{GPS25} & SRS, pairings & $O(1)$ & $O(N)$ & univariate KZG \\
\midrule
Brakedown~\cite{GLSTW23} & transparent, hashes & $O_\lambda(\sqrt N)$ & $O(N)$ field ops & linear-time code \\
BaseFold~\cite{ZCF24} & transparent, hashes & $O_\lambda(\log^2 N)$ & $O(N \log N)$ & sumcheck + folding \\
DeepFold~\cite{GLHQTZ25} & transparent, hashes & $O_\lambda(\log^2 N)$ & $O(N \log N)$ & sumcheck + folding \\
FRI-Binius~\cite{DP24} & transparent, hashes & polylog & $O(N \log N)$ & sumcheck + folding \\
WHIR~\cite{ACFY24} & transparent, hashes & polylog & $\tilde O(N)$ & sumcheck + constr.\ RS \\
KBFold~\cite{Mkh26} & transparent, hashes & $O_\lambda(\log^2 N)$ & $O(N \log N)$ & sumcheck + kernel folding \\
Kronecker-FRI & transparent, hashes & $O_\lambda(\log^2 N)$ & $O(N \log N)$ & extraction + FRI \\
\bottomrule
\end{tabular}
\caption{Multilinear polynomial commitment schemes; SRS is a universal structured reference string (trusted setup). Proof sizes count group elements for the pairing-based schemes, and field elements and hashes for the hash-based ones, with $O_\lambda$ hiding the dependence on the security parameter through the number of queries. The entries for other schemes are as stated in the cited papers.}
\label{tab:apps-compare}
\end{table}

\Cref{tab:apps-compare} lists representative multilinear polynomial commitments.
We do not report timings for schemes that we have not run ourselves: absolute timings from different papers are measured on different machines, with different fields, codes and security targets, and are not comparable.

\paragraph{Pairing-based schemes.}
Gemini, Zeromorph, Mercury and Samaritan reduce a multilinear evaluation to univariate KZG commitments, and reach logarithmic or constant proof size.
Mercury is the closest in spirit to Kronecker-FRI: it also expresses $f(z)$ as a coefficient of a product with a tensor-structured polynomial (\cref{sec:kernel:related}).
These schemes need a universal structured reference string and pairing-friendly curves, and are not post-quantum. Kronecker-FRI is transparent, uses only a hash function, and has the post-quantum guarantee of \cref{cor:pq}, at the price of proofs of about $1$--$2$\,MiB instead of at most a few kilobytes.

\paragraph{Hash-based schemes.}
Among hash-based schemes, Brakedown has a linear-time prover and proofs of size $O(\sqrt N)$; the folding-based schemes have polylogarithmic proofs.
The closest comparison is with the schemes that combine a sumcheck with Reed--Solomon folding, which we measured at identical engineering in \cref{sec:exp}: Kronecker-FRI commits and verifies at the same speed as KBFold and the coefficient-form (BaseFold-style) baseline, has proofs $8$--$17\%$ larger, and opens $3.0$--$3.7$ times more slowly.
DeepFold and WHIR reduce proof size and verification time substantially below these schemes by working in the list-decoding regime, with fewer queries; WHIR reports $63$\,KiB proofs and $360\,\mu$s verification for degree $2^{22}$ at $100$ bits~\cite{ACFY24}; the main theorem of~\cite{ACFY24} (their Theorem~1) assumes a conjecture on correlated agreement (their Conjecture~1), and the authors note that parameters provable without it are less efficient.
Part of the difference is the number of queries: our analysis, like that of KBFold, is confined to the unique-decoding regime (\cref{rem:unique-only}), which forces $\kappa = 148$ at $\rho = 1/4$; the other part is that WHIR's proximity test is itself more query-efficient than FRI~\cite{ACFY24}.

\paragraph{What Kronecker-FRI offers.}
Its advantages are structural rather than in raw speed.
\begin{itemize}
  \item The verifier is an FRI verifier plus an $O(n)$ local computation per query; there is no sumcheck, and no round of the protocol depends on the multilinear structure except the first.
  \item The soundness proof has one algebraic step, the batching lemma, on top of the folding test (\cref{sec:sound}); the folding test itself is independent of the evaluation claim.
  \item The scheme is built from components that already exist in FRI-based systems (\cref{sec:apps:fit}).
  \item Its interactive soundness is fully proved in the unique-decoding regime, without conjectures, from the correlated-agreement theorem of~\cite{BCIKS23}; its post-quantum guarantee is proved for one opening from the quoted theorem of~\cite{CDHZ25} (\cref{cor:pq}), for their compiler.
\end{itemize}
Its costs are an opening $3$--$3.7$ times slower than KBFold's at the same engineering, and the proof sizes of unique-decoding FRI.
The natural next step is a list-decoding analysis of the batching lemma, which would reduce $\kappa$ and bring proof sizes closer to those of DeepFold and WHIR.
